\section{De la escalabilidad en tiempo de prueba a la escalabilidad en tiempo de entrenamiento}

Las clases anteriores mejoraron la tasa de éxito con parámetros fijos mediante muestreo, búsqueda y un verifier. Esta clase pregunta un paso más allá: \textbf{¿se puede convertir esas trayectorias de razonamiento (reasoning trajectory) exitosas en una señal de entrenamiento, de modo que la distribución de una sola salida del modelo mejore por sí misma?}

AIME es el benchmark que motiva principalmente esta clase. Contiene problemas de matemáticas de competencias recientes, lo que reduce el problema de que los benchmarks antiguos estén contaminados por los datos de entrenamiento, y exige más de las cadenas de razonamiento largas, el retroceso y la autoverificación.

\videofigure{figures/aime.jpg}{AIME 2024/2025 se usa para medir la capacidad reciente de razonamiento matemático.}{00:01:16--00:03:28}

El cómputo en tiempo de entrenamiento y el cómputo en tiempo de prueba son dos ejes que se pueden escalar a la vez: el entrenamiento hace que las trayectorias correctas sean más frecuentes en la distribución del modelo; la escalabilidad en tiempo de prueba realiza una búsqueda más profunda o más amplia a partir de esa distribución.

\videofigure{figures/train-test-scaling.jpg}{Aumentar tanto el cómputo en train-time como en test-time mejora el pass@1 en AIME.}{00:05:48--00:07:20}

\begin{importantbox}{El objetivo de la escalabilidad en tiempo de entrenamiento es remodelar la distribución de probabilidad}
El muestreo repetido (repeated sampling) mejora principalmente la probabilidad de «obtener al menos una respuesta correcta»; el RL, en cambio, intenta aumentar en conjunto la probabilidad de las trayectorias con recompensa alta, de modo que la respuesta correcta aparezca con más frecuencia y más consistencia. El resultado ideal es obtener la misma exactitud con menos muestreo en tiempo de prueba.
\end{importantbox}

\subsection{Por qué el RL de razonamiento falla con facilidad}

En una CoT larga, el espacio de acciones (action space) es enorme, la recompensa es escasa y la longitud de las trayectorias varía mucho. El entrenamiento también se enfrenta a problemas de sistema como el reward hacking, el colapso de entropía, la varianza del gradiente, la presión de memoria y el rendimiento del muestreo en línea. Por ello, el curso enfatiza que la «magia» de escalar el RL suele venir de los detalles de implementación, no de una sola función objetivo.

\subsection{Resumen de la sección}

La escalabilidad en tiempo de entrenamiento convierte la recompensa del verifier o del entorno en actualizaciones de parámetros. Puede mejorar la estabilidad de una sola respuesta y del voto mayoritario, pero también puede reducir la diversidad de exploración; entender esta tensión es el hilo común de STaR, GRPO y DAPO.
